\question[4]
Beantworte die folgenden Fragen zum Einerkomplement.
\begin{parts}
\part Gib die Zahl -23 in der 8-Bit Einerkomplement-Darstellung an.
\part Welche Zahl hat in der 8-Bit Einerkomplement-Darstellung die Codierung \texttt{11110010}?
\part Warum gibt es im Einerkomplement zwei Darstellungen der Null? Gib beide für 8 Bit an.
\end{parts}
\begin{solutionbox}{3cm}
a) 23 ist \texttt{00010111}, alle Bits invertiert ergibt \texttt{11101000}. \\
b) Das Bitmuster beginnt mit 1, die Zahl ist also negativ. Invertiert ergibt sich \texttt{00001101} = 13, die Zahl ist -13. \\
c) Invertieren von \texttt{00000000} (+0) ergibt \texttt{11111111} (-0). Beide Muster stehen für die Null.
\end{solutionbox}
